\scanpage{24}

Le troisième pt se ramène au premier, si on introduit pour tt $x\in X$ la ``catégorie fibrée localisée'' $\mathcal F_{(x)}$ par la condition que pour tt ouvert \deleted{$\unsure{W}\ \unsure{w}$} \deleted{\unsure{de}} $w$ de $X_x$, on pose

\[
\mathcal F_{(x)}(\unsure{w|W})=\underset{\longrightarrow}{\mathrm{Pseudolim}}\ \mathcal F(W)
\]

où $W$ parcourt les ouverts de $X$ tels que $W_{(x)}\supset\deleted{\unsure{X_x|W}}w$ (ou si on veut, tels que $W_{(x)}=w$). \deleted{En \unsure{particulier}}

On \unsure{notera|note} que avec cette définition, pour tt ouvert $V$ de $X$ et \deleted{$\unsure{\xi,\eta\in\mathrm{Ob}}$} $\unsure{\xi|\zeta|f},\eta\in\mathrm{Ob}\,\mathcal F(\unsure{U|V})$, on aura \inserted{\deleted{\unsure{évidemment|fonctoriellement}}} un isom can

\[
\underline{\mathrm{Hom}}(\xi,\eta)_{(x)}\xrightarrow{\sim}
\underline{\mathrm{Hom}}(\xi_{(x)},\eta_{(x)})
% double rule below and to the right of the closing parenthesis
\]

comme il résulte aussitôt des définitions.

% slanted
\margin{On a une \uline{flèche} \deleted{\unsure{surj.}} \unsure{triviale}, et on constate que c'est un isom \deleted{\unsure{par}} \unsure{sur} les fibres en tt pt $y\in\unsure{X_{(x)}|V_{(x)}}$}

\deleted{\unsure{Dire|N.B.}} La propriété $P_1$ ($P_2,P_3$) de $\mathcal F$ en $x$, \uline{s'énonce aussi en disant que} $\mathcal F_{(x)}$ \uline{sur} $X_{(x)}$ \uline{satisfait à la propriété} $P_1$ ($P_2,P_3$).

En ce qui concerne $P_1$ $(P_2)$ cela signifie aussi que pour tt \deleted{ouv.} \inserted{vois.} $V\deleted{\subset X}$ \inserted{de $x$}, et $\xi,\eta\in\mathrm{Ob}\,\mathcal F(U)$ $\underline{\mathrm{Hom}}(\xi,\eta)$ satisfait à $P_1$ $(P_2)$ en $x$.

% archive sheet number: 23
